\scanpage{226}
matrices \unsure{$\{$} $\underline u$ \unsure{$\in B$} such that $u\equiv1\ (2)$ i.e. $M\equiv0\ \unsure{(2)}\unsure{\}}$, by the \inserted{closed} subgroup $\underline L_0$ generated by $\ell_0=\begin{pmatrix}1&-2\\0&1\end{pmatrix}$, is identified with the group of matrices $\begin{pmatrix}\mu&0\\0&1\end{pmatrix}$ --- \inserted{\unsure{the}} maximal torus $T_0$ of $B$, more precisely we have a $\frac12$ direct product structure
\[
\begin{array}{ccccc}
\unsure{\underline B_1}&\simeq&T_0&\cdot&\underline L_0\\
\Vert&&\Vert&&\Vert\\
\left\{\begin{pmatrix}\mu&M\\0&1\end{pmatrix}\middle|\begin{array}{l}\mu\in\widehat{\mathbb Z}^{*}\\M\in2\widehat{\mathbb Z}\end{array}\right\}&&
\left\{\begin{pmatrix}\mu&0\\0&1\end{pmatrix}\middle|\mu\in\widehat{\mathbb Z}^{*}\right\}&&
\left\{\begin{array}{c}\ell_0^n\\\Vert\\\begin{pmatrix}1&-2n\\0&1\end{pmatrix}\end{array}\middle|n\in\widehat{\mathbb Z}\right\}
\end{array}
\tag{178}
\]

This shows that if we are interested in the quotient $\mathcal M_{0,3}^{!\sim}(0)/L_0\xrightarrow{\sim}\mathcal N_{0,3}^{!\sim}$, then the only abelian invariant we know how to associate with it in the present context (when we are not working with $\alpha,\beta$) is the multiplier $\mu$. But better still, we find \unsure{thanks} to the decomposition (178) in $\mathrm{GL}(2,\widehat{\mathbb Z})$ a canonical splitting of the extension $\mathcal M_{0,3}^{!\sim}(0)$ of $\mathcal N_{0,3}^{!\sim}$ by $L_0$ :

\uline{Proposition} For $\forall\,\dot u\in\mathcal N_{0,3}^{!\sim}$, there is a unique representative \unsure{$u=u_{\dot u}$|$u=u_{\xi}$} $\in\mathcal M_{0,3}^{!\sim}[0]$ such that we have $\underline u\in T_0$ i.e. $\underline u$ of the form $\begin{pmatrix}\mu&0\\0&1\end{pmatrix}$
\inserted{[i.e. $\xi=\begin{pmatrix}A&B\\C&D\end{pmatrix}$ of the form $\begin{pmatrix}A&0\\C&D\end{pmatrix}=\begin{pmatrix}\unsure{\ast|\frac{1}{D}}&0\\1-D&D\end{pmatrix}$]}, --- In other words, the subgroup

\margin{225}

\scanpage{227}
\mspage{413}
\deleted{\unsure{$\mathcal M_{0,3}^{!\sim}(0)$}}
\[
\mathcal M_{0,3}^{!\sim}(0)\overset{\mathrm{df}}=
\left\{u\in\mathcal M_{0,3}^{!\sim}(0)\ \middle|\ \underline u\in T_0\quad\text{i.e. }\underline u=\begin{pmatrix}\mu&0\\0&1\end{pmatrix}\right\}
\tag{179}
\]
is a section subgroup of the extension $\mathcal M_{0,3}^{!\sim}[0]$ of $\mathcal N_{0,3}^{!\sim}$ by $L_0$, hence also of the extension $\mathcal M_{0,3}^{!\sim}$ of $\mathcal N_{0,3}^{!\sim}$ by $\widehat\Pi_{0,3}$, \unsure{or} of the extension $\mathcal M_{0,3}^{\sim}$ of $\mathcal N_{0,3}^{!\sim}$ by $\widehat{\mathcal G}_{0,3}^{+}$.

\unsure{This|That} comes from the fact that if $u\in\mathcal M_{0,3}^{!\sim}[0]$ \inserted{defines} $\underline u=\begin{pmatrix}\mu&M\\0&1\end{pmatrix}$, thus with $M\equiv0\ (2)$, i.e. $M\in2\widehat{\mathbb Z}$, there is a \uline{unique} $n\in\widehat{\mathbb Z}$ such that $\underline{\ell_0^n u}=\begin{pmatrix}\mu&M-2n\\0&1\end{pmatrix}$ is in $T_0$ i.e. such that $\ell_0^n u\in\mathcal M_{0,3}^{!\sim}(0)$, namely $n=\dfrac M2$.

If $(\xi,\deleted{\unsure{\varepsilon}}\eta)$ satisfy conditions (133)(134) i.e. define a $u=u_{\xi,\eta}$ (cf. Theorem page~407'), then we shall say that $(\xi,\eta)$ is \uline{normalized} if we have $u_{\xi,\eta}\in\mathcal M_{0,3}^{!\sim}(0)$. Likewise, if an arbitrary element of $\mathcal M_{0,3}^{!\sim}$ is written in the form
\[
u=\cdot f^{-1}\cdot u_{\xi,\eta}=f^{-1}u_{\unsure{\xi}},
\]
we say that this representation \inserted{(\unsure{tainted} with the indeterminacy (136))} is \uline{normalized}, if $u_{\xi,\eta}=u_\xi$ is normalized, [i.e. \unsure{if} \inserted{for} $\xi\unsure{\simeq}\begin{pmatrix}A&B\\C&D\end{pmatrix}$ \deleted{the form} we have $B=0$]. Thus, \deleted{every} $u$ can be written uniquely in the form
\[
u_{f,\xi,\eta}=f^{-1}u_\xi,
\tag{\unsure{179}}
\]
\margin{226}

\scanpage{228}
with $\unsure{f|\xi}$ normalized. In fact, this corresponds to the $\frac12$-direct product decomposition
\[
\mathcal M_{0,3}^{!\sim}\simeq\mathcal M_{0,3}^{!\sim}(0)\cdot\widehat\Pi_{0,3},
\tag{180}
\]
moreover we likewise have
\[
\mathcal M_{0,3}^{\sim}\simeq\mathcal M_{0,3}^{!\sim}(0)\cdot\widehat{\mathcal G}_{0,3}^{+}.
\tag{181}
\]

\uline{Remark} We should at some point \unsure{try} to relate these canonical splittings to the choice of a uniformizer at $0$ of $\mathbb P^1_{\mathbb Q}=\widehat U_{0,3\mathbb Q}$. Of course, the same question also arises \unsure{at|for} the \inserted{other two} ``points at infinity'' $1,\infty$ of $U_{0,3\mathbb Q}$, relative to the section subgroups
\[
\mathcal M_{0,3}^{!\sim}(q)\qquad q\in\{0,1,\infty\}
\tag{182}
\]
corresponding to these points at infinity --- the construction of a section-group $\mathcal M_{0,3}^{!\sim}(0)$ in the normalizer of $L_0$, obviously applies without change for $L_1,L_\infty$. We shall moreover have
\[
\rho\bigl(\mathcal M_{0,3}^{!\sim}(q)\bigr)=\mathcal M_{0,3}^{!\sim}\bigl(\unsure{\dot\rho|\rho}(q)\bigr)\qquad\text{for }q\in\{0,1,\infty\},
\tag{183}
\]
so the \inserted{three} $\mathcal M_{0,3}^{!\sim}(q)$ are obtained from any one of them by applying $\rho,\rho^2$.
\margin{227}

\scanpage{229}
\mspage{415}
47. Return to the homomorphism $\widehat\Pi_{0,3}\xrightarrow{\widehat\varphi}\widehat{\mathcal G}_{0,3}^{+}$ and the condition of commutation with \uline{it} up to inner automorphisms. The groups $\mathcal M_{0,3}^{!\approx},\mathcal M_{0,3}^{\approx}$ and the homomorphism $\Phi:\mathcal M_{0,3}^{!\approx}\longrightarrow\mathcal M_{0,3}^{\approx}$.

Recall \inserted{(\S41)} that $\widehat\varphi$ is described by
\[
\widehat\varphi(\ell_0)=\varepsilon_0=\sigma_\infty\rho,\quad
\widehat\varphi(\ell_1)=\sigma_\infty,\quad
\widehat\varphi(\ell_\infty)=\rho^{-1}.
\tag{1}
\]
Let
\[
u=u_{\alpha,\beta,f}\quad\in\mathcal M_{0,3}^{!\sim}
\]
be a loop automorphism of $\widehat\Pi_{0,3}$, commuting with $\mathcal G_3$ up to inner automorphisms, defined by
\[
\begin{aligned}
u(\sigma_\infty)&=h\sigma_\infty=\operatorname{int}(\beta_0)\sigma_\infty
&\left(\begin{array}{l}h=\beta_0\sigma_\infty(\beta_0)^{-1},\\\beta_0=f^{-1}\beta\end{array}\right)
\end{aligned}
\tag{2}
\]
\[
\begin{aligned}
u(\rho)&=g\rho=\operatorname{int}(\alpha_0)\rho
&\left(\begin{array}{l}g=\alpha_0\rho(\alpha_0)^{-1},\ \deleted{\unsure{=}}\\\alpha_0=f^{-1}\alpha\end{array}\right)
\end{aligned}
\tag{3}
\]
and let us write that $u$ commutes with $\varphi$ up to inner automorphisms, i.e. that there exists \deleted{\unsure{$\lambda\in$}}
\[
\lambda\in\widehat\Pi_{0,3}
\tag{4}
\]
such that we have
\[
\widehat\varphi u(x)=\operatorname{int}(\lambda)u\widehat\varphi(x)
\qquad\forall x\in\widehat\Pi_{0,3}.
\tag{5}
\]
It suffices to check (5) for $x=\ell_0$ and $x=\ell_1$.

1) Case $x=\ell_0$ \deleted{\unsure{$=(\sigma_\infty\rho)^2$}}, whence
\[
u(x)=\deleted{u(\sigma_\infty)u(\rho)}\deleted{\unsure{=h\sigma_\infty g\rho=h(\sigma_\infty g\sigma_\infty^{-1})\sigma_\infty\rho}}
\inserted{\operatorname{int}(f^{-1})(\ell_0^\mu)}
\]
\[
\widehat\varphi(u(x))=\operatorname{int}\bigl(\widehat\varphi(f)^{-1}\bigr)\widehat\varphi(\ell_0^\mu)=\operatorname{int}\bigl(\widehat\varphi(f)^{-1}\bigr)\varepsilon_0^\mu
\]
\[
u\widehat\varphi(x)=u(\varepsilon_0)=\operatorname{int}(f^{-1})(\varepsilon_0^\mu)
\]
so the relation (5) reads
\[
\operatorname{int}\bigl(\deleted{\unsure{f\lambda^{-1}}}\widehat\varphi(f)\unsure{\lambda}f^{-1}\bigr)(\varepsilon_0^\mu)=\varepsilon_0^\mu
\]
\margin{228}

\scanpage{230}
is equivalent to
\[
\deleted{\unsure{f\lambda^{-1}}}\widehat\varphi(f)\lambda f^{-1}\in L_0^{\sim}\quad
\left(=\operatorname{Centr}_{\widehat\Pi_{0,3}}(\varepsilon_0^\mu)\right)
\]
\margin{N.B. we do not necessarily have $\widehat\varphi(f)\in\widehat\Pi_{0,3}$, so \unsure{neither} \illegible{1} $\widehat\varphi(f)\lambda f^{-1}\in\widehat\Pi_{0,3}\ldots$}
i.e. to the existence of $r\in\widehat{\mathbb Z}$ such that we have
\[
\deleted{\unsure{f\lambda^{-1}}}\widehat\varphi(f)\lambda f^{-1}=\varepsilon_0^r,
\]
i.e.
\[
\lambda=\widehat\varphi(f)^{-1}\varepsilon_0^r f.
\tag{6}
\]
This thus makes it possible to eliminate $\unsure{\exists\lambda}$ from the second relation :

2) Case $x=\ell_1$, whence
\[
u(x)=\operatorname{int}(f_1^{-1})(\ell_1^\mu)
\qquad\bigl(f_1^{-1}=\deleted{\unsure{\rho(f)}}g\rho(f^{-1})\bigr)
\]
\[
\widehat\varphi(u(x))=\operatorname{int}\bigl(\widehat\varphi(f_1)^{-1}\bigr)\sigma_\infty^\mu=\operatorname{int}\bigl(\widehat\varphi(f_1)\bigr)\sigma_\infty
\]
(taking into account that $\sigma_\infty^\mu=\sigma_\infty$ since $\mu\equiv1\ (2)$).
\deleted{\unsure{$u\widehat\varphi(x)=\operatorname{int}(\widehat\varphi(f_1))$}}
\[
u\widehat\varphi(x)=u(\sigma_\infty)=\operatorname{int}(\beta_0)(\sigma_\infty),
\]
So the relation (5) reads
\[
\operatorname{int}\bigl(\widehat\varphi(f_1)^{+1}\lambda\beta_0\bigr)\sigma_\infty=\sigma_\infty,
\]
which also means that there exists $\unsure{i|\alpha}\in\mathbb Z/2\mathbb Z$ such that we have \deleted{\illegible{1}}
\[
\widehat\varphi(f_1)^{+1}\lambda\beta_0=\sigma_\infty^i
\quad\text{i.e.}\quad
\widehat\varphi(f_1^{-1})^{+1}=\lambda\beta_0\sigma_\infty^i,
\]
and replacing $\lambda$ by its value (6), and $f_1^{-1}$ by its value $g\rho(f^{-1})\deleted{\unsure{f}}$, we find the condition
\[
\deleted{\unsure{\widehat\varphi(\rho(f)^{-1})\widehat\varphi(g)^{-1}=\widehat\varphi(f)^{-1}\ell_0^r f\beta_0\sigma_\infty^\alpha}}
\]
\[
\widehat\varphi\bigl(g\rho(f)^{-1}\bigr)=\widehat\varphi(f)^{-1}\varepsilon_0^r\underbrace{f\beta_0}_{\beta}\sigma_\infty^{\inserted{\unsure{\alpha}}}\quad\text{i.e.}
\]
\[
\widehat\varphi\bigl(\underbrace{fg\rho(f)^{-1}}_{\xi=\alpha\rho(\alpha)^{-1}}\bigr)=\varepsilon_0^r\beta\sigma_\infty^\alpha
\]
\margin{229}

\scanpage{231}
\mspage{416}

Thus :

\uline{\deleted{Proposition}\inserted{Theorem}} For
\[
\widehat\varphi:\widehat\Pi_{0,3}\longrightarrow
\widehat{\unsure{\mathcal G}}_{0,3}^{+}
\]
to commute up to inner automorphisms with the element
$u\in\mathcal M_{0,3}^{!\sim}$, corresponding to the invariants
$g,h,f,\alpha_0,\beta_0,\alpha,\beta,\xi,\eta$ (cf. \S 46),
it is necessary and sufficient that we have
\[
\widehat\varphi(\xi)=\varepsilon_0^r\beta\sigma_\infty^{\unsure{i|\alpha}}
\qquad\text{for }r\in\widehat{\mathbb Z},\quad
\alpha\in\mathbb Z/2\mathbb Z\ \text{suitable}.
\tag{7}
\]
\margin{\uline{N.B.} $r,\alpha$ are well determined when $\alpha$ and $\beta$
(or again $\xi$ and $\beta$) are chosen.}
We shall then have (5), with \deleted{Note that}
$\lambda$ given by (6) ;
$\lambda=\widehat\varphi(f)^{-1}\ell_0^r f$.

Note that $u$ defines the triple $(\alpha,\beta,f)$ only modulo
a transformation of the form
\deleted{$(\alpha,\beta,f)\ \unsure{\mapsto}\ \unsure{\alpha'=\ell_0^n\alpha,\ \beta}$}
\[
(\alpha,\beta,f)\longmapsto
(\alpha'=\ell_0^n\alpha,\ \beta'=\ell_0^n\beta,\ f'=\ell_0^n f)
\tag{8}
\]
where $n\in\widehat{\mathbb Z}$. We shall then have (putting
$\xi'=\alpha'\rho(\alpha')^{-1}=\ell_0^n\xi\ell_1^{-n}$)
\[
\widehat\varphi(\xi')=\varepsilon_0^{r'}\beta'
\sigma_\infty^{\unsure{i'|\alpha'}}
\tag{*}
\]
with $r'\in\widehat{\mathbb Z}$,
$\unsure{i'|\alpha'}\in\mathbb Z/2\mathbb Z$ also well determined.
Now we shall have
\[
\begin{aligned}
\widehat\varphi(\unsure{\xi'})
 &=\widehat\varphi(\ell_0^n\unsure{\xi}\ell_1^{-n})\\
 &=\widehat\varphi(\ell_0^n)\widehat\varphi(\unsure{\xi})
   \widehat\varphi(\ell_1^{-n})\\
 &=\varepsilon_0^n\widehat\varphi(\unsure{\xi})\sigma_\infty^{-n}
\end{aligned}
\]
\deleted{\illegible{1}}\deleted{\illegible{1}}
and the relation (*) reads
\[
\widehat\varphi(\xi)=\varepsilon_0^{\unsure{r'-n|r'-2n}}\beta'
\sigma_\infty^{\unsure{i'-n|\alpha'-n}}
\]
or again, since $\beta'=\ell_0^n\beta=\varepsilon_0^{2n}\beta$
\[
\widehat\varphi(\xi)=\varepsilon_0^{r'+n}\beta
\sigma_\infty^{\unsure{\alpha'-n|i'-n}},
\]
which, compared with (7), gives
\[
r'=r-n,\qquad\unsure{i'|\alpha'}=\unsure{i|\alpha}+n\ \mathrm{mod}\,2.
\tag{9}
\]
\margin{230}

\scanpage{232}

This implies the

\uline{Corollary 1} Let us suppose that $u$ commutes with
$\widehat\varphi:\widehat\Pi_{0,3}\longrightarrow
\widehat{\mathcal G}_{0,3}^{+}$ up to inner automorphisms. Then it is possible
\uline{in a unique way} to describe \unsure{$u$|$\varphi$} by a triple
$(\alpha,\beta,f)$ (a priori determined modulo a transformation of type (8),
in such a way that we have (putting $\xi=\alpha\rho(\alpha)^{-1}$,
as usual) \unsure{a} relation
\[
\boxed{\widehat\varphi(\xi)=\beta\sigma_\infty^i}
\qquad\text{for }i\in\mathbb Z/2\mathbb Z\ \text{suitable}.
\tag{10}
\]
\margin{We should examine whether this definition coincides with that on
\unsure{the} page\,\unsure{410'}\,,\,\unsure{413|417}, formulated without
involving $\widehat\varphi$.}
Such a choice \unsure{will be} called \uline{normalized}.

Thus the $u\in\mathcal M_{0,3}^{!\sim}$ which commute with
$\widehat\varphi$ up to inner automorphisms correspond one-to-one to the
\deleted{pairs} triples \deleted{$(\alpha,\beta,f)$}
\deleted{satisfying (10)}
\[
(\alpha,\beta,f),\qquad
\alpha\in\unsure{\widetilde{\widehat\Pi}_{0,3}},\quad
\beta,f\in\widehat\Pi_{0,3}
\]
satisfying, on the one hand the loop relation
\[
\xi\unsure{\doteq|=}\eta\ell_1^\nu
\qquad
\begin{gathered}
\text{where }\xi=\alpha\rho(\alpha)^{-1},\quad
\eta=\beta\sigma_\infty(\beta)^{-1},\\
\text{and where }\nu\in\widehat{\mathbb Z}\ \text{suitable}\\
(\text{such that }2\nu+1\in\widehat{\mathbb Z}^{*})
\end{gathered}
\tag{11}
\]
on the other hand relation (10) (for suitable $i\in\mathbb Z/2\mathbb Z$)
--- neither relation moreover involving $f$.
\margin{231}

\scanpage{233}
\mspage{417}

Let us note that $i\in\mathbb Z/2\mathbb Z$ is defined in terms of
$\xi$ (taking into account
$\beta\in\widehat\Pi_{0,3}=\operatorname{Ker}
(\widehat{\mathcal G}_{0,3}^{+}\to\unsure{\mathcal G_3|\mathfrak S_3})$)
by reducing (10) modulo $\widehat\Pi_{0,3}$. Denoting by
\[
\widehat\varphi^{\bullet}:
\widehat\Pi_{0,3}\longrightarrow\unsure{\mathcal G_3|\mathfrak S_3}
\tag{12}
\]
the composite
\[
\widehat\Pi_{0,3}\xrightarrow{\widehat\varphi}
\widehat{\mathcal G}_{0,3}^{+}
\overset{\mathrm{can}}{\longrightarrow}
\unsure{\mathcal G_3|\mathfrak S_3},
\]
formula (10) also gives
\[
\widehat\varphi^{\unsure{\bullet|\circ}}(\xi)=\sigma_\infty^i
\tag{13}
\]
which does determine $i$\inserted{$\in\mathbb Z/2\mathbb Z$} in terms of
$\xi\in\widehat\Pi_{0,3}$, and implies an intrinsic condition on
the element $\xi$ of $\widehat\Pi_{0,3}$, namely that its image in
$\unsure{\mathcal G_3|\mathfrak S_3}$ under
$\widehat\varphi^{\unsure{\bullet|\circ}}$ should indeed be in the subgroup
$\{1,\sigma_\infty\}$ of $\unsure{\mathcal G_3|\mathfrak S_3}$.
In fact, the kernel of $\widehat\varphi^{\unsure{\bullet|\circ}}$ is
the \unsure{inverse} subgroup of $\widehat\Pi_{0,3}$ generated by
$\ell_0^2,\ell_1^2,\ell_\infty^3$, and the condition being considered
therefore means that modulo this subgroup, $\xi$ should be of the form $\ell_1^i$.

By formula (10), we can eliminate $\beta$ from the datum
$(\alpha,\beta,f)$, and we can replace the datum $\alpha$ by the datum
$\xi$. Thus:
\margin{232}

\scanpage{234}

\uline{Corollary 2} \unsure{The} elements of the \unsure{$^{s}/$group} of
$\mathcal M_{0,3}^{!\sim}$ formed by the $u$ which commute with
$\widehat\varphi$ up to inner automorphisms, is \deleted{\illegible{1}} in one-to-one correspondence with
the set of pairs $(\xi,f)$ \deleted{\unsure{$\in$}}
$(\xi,f\in\widehat\Pi_{0,3})$ satisfying the conditions
(\unsure{not|by} involving anything but $\xi$)
\[
\begin{array}{c@{\qquad}c}
\begin{gathered}(14)\\[1.8ex](15)\end{gathered}
&
\boxed{\begin{gathered}
\xi\rho(\xi)\rho^2(\xi)=1\\[1.8ex]
\xi=\bigl(\widehat\varphi(\xi)\,
\sigma_\infty(\widehat\varphi(\xi))^{-1}\bigr)\ell_1^\nu
\end{gathered}}
\end{array}
\]
\[
(\nu\in\widehat{\mathbb Z}\ \text{suitable such that }
2\nu+1\in\widehat{\mathbb Z}^{*}).
\]
\margin{we also need the ``preliminary'' ``condition''\ldots}

For such a $\xi$, $\widehat\varphi^{\unsure{\bullet|\circ}}(\xi)$
in $\unsure{\mathcal G_3|\mathfrak S_3}$ is of the form
$\sigma_\infty^i$, with $i\in\mathbb Z/2\mathbb Z$, and putting
\[
\left\{\begin{aligned}
&\beta=\widehat\varphi(\xi)\sigma_\infty^i
  \quad(\in\widehat\Pi_{0,3})\\
&\alpha\in\widehat\Pi_{0,3}\quad
  \text{the unique element of }\widetilde{\widehat\Pi}_{0,3}
  =\widehat\Pi_{0,3}\cdot\{1,\tau_\infty\}\\
&\hspace{5em}\text{such that }\alpha\rho(\alpha)^{-1}=\xi
\end{aligned}\right.
\tag{16}
\]
\margin{in fact, we do not really need $\alpha,\beta$}
$u$ is obtained from $\xi,f$ by the formulas
\[
\left\{\begin{aligned}
u(\sigma_\infty)
 &=\operatorname{int}(f^{-1}\unsure{\beta})(\sigma_\infty)\\
 &=\operatorname{int}(f^{-1})(\operatorname{int}(\unsure{\beta})(\sigma_\infty))\\
 &=\operatorname{int}(f^{-1})(\xi\ell_1^{-\nu}\sigma_\infty)\\
u(\rho)
 &=\operatorname{int}(f^{-1}\alpha)(\rho)\\
 &=\operatorname{int}(f^{-1})(\operatorname{int}(\alpha)(\rho))\\
 &=\operatorname{int}(f^{-1})(\unsure{\xi}\rho)
\end{aligned}\right.
\tag{17}
\]

We shall note that in these formulas, $f$ enters only through
composition with $\operatorname{int}(f)$. If we
\deleted{put}\inserted{denote}, for $\xi\in\widehat\Pi_{0,3}$
\margin{233}

\scanpage{235}
\mspage{418}

\margin{This \unsure{notation} is already in use for the
general \unsure{considerations} of the preceding \S
$\unsure{u_{\xi,\eta}}$ depends in \unsure{any} case only on $\xi$\ldots}
satisfying (14)(15), by $u_\xi$ the automorphism
\deleted{$u_\xi(\sigma_\infty)$} corresponding, in terms of
Corollary 2, to the pair $(\xi,1)$, so that
\[
\left\{\begin{aligned}
u_\xi(\sigma_\infty)
 &=\deleted{\unsure{\operatorname{int}(\alpha)}}
 \overbrace{(\xi\ell_1^{-\nu})}^{\unsure{\eta|\nu}}\sigma_\infty\\
u_\xi(\rho)&=\xi\rho
\end{aligned}\right.,
\tag{18}
\]
formulas (17) can also be written, simply
\[
u_{\xi,f}=\operatorname{int}(f^{-1})\,u_\xi.
\tag{19}
\]

On the other hand formulas (5), (6) give here
\[
\widehat\varphi\circ u_\xi=u_\xi\circ\widehat\varphi
\tag{20}
\]
where, of course, on the first \unsure{side|side} $u_\xi$ denotes the
restriction of $u_\xi$ to $\widehat\Pi_{0,3}$.

Let us note

\uline{Corollary 3} The loop automorphisms of
$\widehat{\unsure{\mathcal G}}_{0,3}^{+}$ of the form $u_\xi$
\inserted{($\xi$ \unsure{normalized|normalized}$^*$)} are exactly those for
which we have the \deleted{properties} formulas (20) (i.e. which
\deleted{\illegible{1}} commute with $\widehat\varphi$) and which moreover
normalize $L_0$.
\margin{$^*$ in the sense of page\,416'}
\figure{\begin{tikzpicture}[x=1mm,y=1mm,line width=.3pt,line cap=round]
\draw (0,35) .. controls (-.16,29) and (.15,23) .. (-.02,17);
\draw (.65,35.0) -- (.88,34.1) -- (1.38,34.55);
\draw (0,9.5) .. controls (-.08,6.3) and (-.16,2.7) .. (0,0);
\end{tikzpicture}}
\margin{234}

\scanpage{236}

Indeed, in terms of the \inserted{triples} $(\alpha,\beta,f)$ \inserted{``special''} these properties are expressed by the \inserted{\unsure{simple}} conditions \deleted{that we have} $f$ of the form $\ell_0^n$ (which expresses the condition of normalizing $L_0$) and that $\lambda$ given by (6) \unsure{case|for} $r=0$, i.e.
\[
\lambda=\widehat\varphi(f)^{-1}f
\tag{21}
\]
be equal to $1$, i.e. the condition
\[
\widehat\varphi(f)=f,
\tag{22}
\]
which in this case reads
\[
\varepsilon_0^n=\underbrace{\ell_0^n}_{\varepsilon_0^{2n}}
\qquad\text{i.e.}\quad \varepsilon_0^n=1
\]
which implies $\unsure{f|f_0}=\ell_0^n=1$, hence $u$ of the form $u_{(\alpha,\beta,1)}=u_\xi$, where $\xi=\alpha\rho(\alpha)^{-1}$.

\uline{Corollary 4} The $u_\xi$ \inserted{normalized} form a \uline{subgroup} of $\operatorname{Norm}_{\mathcal M_{0,3}^{!\sim}}(L_0)$, i.e. $\mathcal M_{0,3}^{!\approx}(0)$, whose \deleted{\unsure{image}}\inserted{intersection} with $\widehat\Pi_{0,3}$ is reduced to $\{1\}$, and whose image in $\mathcal N_{0,3}^{!\sim}$ is \inserted{the subgroup $\mathcal N_{0,3}^{!\approx}$} formed by the outer loop automorphisms of $\widehat{\mathcal G}_{0,3}^{+}$, stabilizing $\widehat\Pi_{0,3}$ (i.e. $u\in\unsure{\mathcal M_{0,3}^{+\widehat{\ }}}$) such that $\widehat\varphi:\widehat\Pi_{0,3}\to\widehat{\mathcal G}_{0,3}^{+}$ commutes with $u$ up to inner automorphisms. We therefore have an isomorphism
\[
\mathcal M_{0,3}^{!\approx}(\unsure{0})
=\{u_\xi\mid \xi\in\widehat\Pi_{0,3}\text{ solution of }(14)(15)\}
\xrightarrow{\sim}\mathcal N_{0,3}^{!\approx}.
\tag{\unsure{23}}
\]
\margin{235}

\scanpage{237}
\mspage{419}

It only remains to check that a $u_\xi$ can be an inner automorphism only if $u_\xi=1$ i.e. $\xi=1$ (hence $\nu=0$). Now $u_\xi$ is a $u_{\alpha,\beta}=u_{\alpha,\beta,1}$ ($\alpha,\beta$ defined \deleted{by} via $\xi$ by $\alpha\rho(\alpha)^{-1}=\xi$, $\beta\sigma_\infty(\beta)^{-1}=\xi \ell_1^{-\nu}$), which is an inner automorphism iff we have $\alpha=\ell_0^n$, $\beta=\ell_0^n$ (cf. proposition on page 404), i.e. $\xi=\ell_0^n \ell_1^{-n}$, $\eta=\beta\sigma_\infty(\beta)^{-1}=\ell_0^n \ell_1^{-n}$, and \deleted{the relation} \deleted{\unsure{for}} $\xi=\eta \ell_1^\nu$ reads $\ell_1^\nu=1$ i.e. $\nu=0$. Relation (15) therefore becomes (since
\[
\begin{gathered}
\widehat\varphi(\xi)=\widehat\varphi(\ell_0^n \ell_1^{-n})=\varepsilon_0^n\sigma_\infty^n\\
\sigma_\infty(\widehat\varphi(\xi))=\varepsilon_1^n\sigma_\infty^n
\end{gathered}
\]
) \deleted{\unsure{gives}} the relation
\[
\ell_0^n \ell_1^{-n}=\varepsilon_0^n\sigma_\infty^{-n}\varepsilon_1^n\sigma_\infty^n,
\]
i.e. (since $\ell_0=\varepsilon_0^2$, $\ell_1=\varepsilon_1^2$) if $n$ even
\[
\varepsilon_0^n=\deleted{\unsure{\sigma_\infty}}\varepsilon_1^{3n}
\qquad\text{whence }n=0,
\]
and if $n$ odd
\[
\varepsilon_0^n=\deleted{\unsure{\varepsilon_0}}\underbrace{\sigma_\infty\varepsilon_1^n\sigma_\infty}_{\sigma_\infty(\varepsilon_1)^n=\varepsilon_0^n}\varepsilon_1^{2n}
\]
i.e. $\varepsilon_1^{2n}=0$, i.e. again $n=0$, q.e.d.

\margin{\uline{N.B.} \illegible{1} $\unsure{\mathcal M_{0,3}^{!\sim}[0]}$ --- \unsure{a} \illegible{2} $\mathcal N_{0,3}^{!\sim}$ in its entirety, \unsure{but it} remains to be seen whether \unsure{the other}\ldots\ \illegible{1}}

\uline{\unsure{Scholium}} We have thus found a canonical partial section of $\mathcal M_{0,3}^{!\sim}$ over the subgroup $\mathcal N_{0,3}^{!\approx}$ of $\mathcal N_{0,3}^{!\sim}$. Since
\margin{236}

\scanpage{238}

this subgroup contains $\Gamma_{\mathbb Q}$, we therefore find a section of the extension $\mathcal M_{0,3}^{!}$ of $\Gamma_{\mathbb Q}$ by $\widehat\Pi_{0,3}$, \deleted{It is necessary} \inserted{which is \unsure{in|in}} fact a section of the subgroup of $\mathcal M_{0,3}^{!}$ which is the normalizer of $L_0$. We should check that this canonical section is associated with the choice of a uniformizer \unsure{at the} point $0$ of $\widehat U_{0,3\mathbb Q}$, \deleted{\unsure{and which will be determined}} \inserted{whose image in} $\unsure{\mathfrak m_0}/\mathfrak m_0^2$ should be \inserted{well} determined, where $\mathfrak m_0$ denotes the maximal ideal of the local ring at $0$ in $\widehat U_{0,3\mathbb Q}=\mathbb P_{\mathbb Q}^1$ \unsure{)} --- we should determine this element of $\mathfrak m_0/\mathfrak m_0^2$ --- We feel like taking an element which is, not only a basis over $\mathbb Q$, but even (in the situation over the integers) a basis over $\mathbb Z$ --- \deleted{\unsure{this}} \inserted{condition} which indeed determines an element of $\mathfrak m_0/\mathfrak m_0^2$ \uuline{up to} \uline{sign} --- \unsure{this would be} in fact the choice corresponding to the uniformizer $z$, or $-z$. But is it really the right one? We should come back to this!
\margin{237}

\scanpage{239}
\mspage{420}

Let us work henceforth in the subgroup $\mathcal M_{0,3}^{!\approx}$ of $\mathcal M_{0,3}^{!\sim}$, consisting of the $u$ which commute as outer actions with $\widehat\varphi$, i.e. inverse image of $\mathcal N_{0,3}^{!\approx}$, so that we have
\[
\diagram{\begin{tikzcd}[ampersand replacement=\&]
1 \arrow[r] \& \widehat\Pi_{0,3} \arrow[r] \& \mathcal M_{0,3}^{!\approx} \arrow[r] \& \underset{\substack{\cup\\\Gamma_{\mathbb Q}}}{\mathcal N_{0,3}^{!\approx}} \arrow[r] \& 1
\end{tikzcd}}
\tag{24}
\]
We have therefore found a section subgroup
\[
\mathcal M_{0,3}^{!\approx}(0)\xrightarrow{\sim}\mathcal N_{0,3}^{!\approx}
\tag{25}
\]
\deleted{$\{\unsure{\xi}\mid\xi\in\widehat\Pi_{0,3}\text{ solution of }(14)(15)\}$ \illegible{2}}
\[
\{u_\xi\mid\xi\in\widehat\Pi_{0,3}\text{ solution of }(14)\text{ and }(15)\}\subset\mathcal M_{0,3}^{!\approx}.
\]

In terms of this we can therefore write $\mathcal M_{0,3}^{!\approx}$ as a $\frac12$ direct product
\[
\mathcal M_{0,3}^{!\approx}\simeq\mathcal M_{0,3}^{!\approx}(\unsure{0})\cdot\widehat\Pi_{0,3},
\tag{26}
\]
by writing every element of $\mathcal M_{0,3}^{!\approx}$ uniquely as
\[
\begin{gathered}
\deleted{\unsure{u_{\xi,f}}}\operatorname{int}(f^{-1})\circ\unsure{u_0|u_\xi}\\
\bigl(\unsure{u_0|u_\xi}\in\mathcal M_{0,3}^{!\approx}(0)\text{ i.e. }u_0=u_\xi\text{ with }\xi\text{ solution of }(14)(15)\bigr)\\
f\in\widehat\Pi_{0,3}\text{ well determined}
\end{gathered}
\tag{27}
\]
These are precisely the \deleted{$u_{\xi,f}=u_{\alpha,\beta,f}$} \inserted{$u_{\xi,f}$ of (19), i.e.,}\margin{\illegible{1}}
\[
u_{\xi,f}=u_{\alpha,\beta,f}
\tag{28}
\]
with $\alpha,\beta$ determined in terms of $\xi$ by
\[
\alpha\rho(\alpha)^{-1}=\xi,\quad\deleted{\unsure{\varepsilon}}\beta\sigma_\infty(\beta)^{-1}=\xi \ell_1^{-\nu},
\tag{29}
\]
\margin{238}

\scanpage{240}

this is how we can interpret the \uline{normalized} \unsure{triples} $(\alpha,\beta,f)$ of corollary 1.

\uline{\unsure{Rescholium}} \inserted{(Recapitulation)} The group $\mathcal N_{0,3}^{!\approx}\simeq\mathcal M_{0,3}^{!\approx}(0)$ is in $1$-$1$ correspondence with \unsure{the set} of $\xi\in\widehat\Pi_{0,3}$ satisfying the two relations
\[
\left\{
\begin{aligned}
\xi\rho(\xi)\rho^2(\xi)&=1\\
\xi&=\bigl(\widehat\varphi(\xi)\sigma_\infty(\widehat\varphi(\xi))^{-1}\bigr)\ell_1^\nu
\end{aligned}
\right.
\qquad\begin{gathered}\text{for}\\\nu\in\widehat{\mathbb Z}\text{ suitable}\end{gathered}
\tag{30}
\]
or again (putting $\xi=\alpha\rho(\alpha)^{-1}$, $\alpha\in\unsure{\widehat{\widetilde\Pi}_{0,3}}$) with \unsure{the set} of $\alpha\in\unsure{\widehat{\widetilde\Pi}_{0,3}}$ satisfying the \inserted{second} relation (30), in which we replace $\xi$ by $\alpha\rho(\alpha)^{-1}$. The element $u_\xi$ of $\mathcal M_{0,3}^{!\approx}(0)$ defined by $\xi$ is defined by the formulas (18), i.e. by
\[
u_\xi(\sigma_\infty)=(\xi \ell_1^{-\nu})\sigma_\infty,
\qquad u_\xi(\rho)=\xi\rho.
\tag{31}
\]
Finally, if we have $u_\xi$ and $u_{\xi'}$, then
\[
u_{\xi'}u_\xi=u_{\xi''},
\tag{32}
\]
where $\xi''$ is given by the formulas on page \unsure{395'} (52) --- noting that here $\xi=g$, $\xi'=g'$, $\xi''=g''$:
\deleted{$\xi''=\alpha''\rho(\alpha'')^{-1}$ with $\alpha''=u_{\xi'}(\alpha)\cdot\alpha'$, \unsure{whence|$f''=1$}}
\[
\deleted{\xi''=\deleted{\unsure{\alpha''\rho}}u_{\xi'}(\alpha)\underbrace{\alpha'\rho(\alpha')^{-1}}_{\xi'}\rho(u_{\xi'}(\alpha))^{-1}}
\]
\margin{239}

\scanpage{241}
\mspage{421}

\deleted{\illegible{1}}
\[
\xi''=u'(\xi)\cdot\xi'.
\tag{33}
\]

Let us now set
\[
\mathcal N_{0,3}^{\approx}
=\mathcal G_3\times\mathcal N_{0,3}^{!\approx}
\subset\mathcal N_{0,3}^{\sim}
=\mathcal G_3\times\mathcal N_{0,3}^{!\sim}
\tag{34}
\]
and let $\mathcal M_{0,3}^{\approx}\subset\deleted{\unsure{\mathcal G}}\mathcal M_{0,3}^{\sim}$ be the inverse image in $\mathcal M_{0,3}^{\sim}$ of $\mathcal N_{0,3}^{\approx}\subset\mathcal N_{0,3}^{\sim}$, so that we have an extension structure
\[
1\longrightarrow\deleted{\unsure{\widehat{\mathcal G}_{0,3}^{+}}}\widehat\Pi_{0,3}
\longrightarrow\mathcal M_{0,3}^{\approx}
\unsure{\twoheadrightarrow|\longrightarrow}
\underset{\simeq\mathcal G_3\times\mathcal N_{0,3}^{!\approx}}{\mathcal N_{0,3}^{\approx}}
\longrightarrow1
\tag{35}
\]
containing the \deleted{extension}\inserted{\unsure{subextension}} (24) as the inverse image of \deleted{\unsure{the}}
\deleted{\illegible{7}}
\[
\mathcal N_{0,3}^{!\approx}
=\operatorname{Ker}(\mathcal N_{0,3}^{\approx}\longrightarrow\mathcal G_3).
\]
\deleted{\unsure{(36)}} \inserted{\unsure{We have}} \deleted{\unsure{also}} a \unsure{further} \unsure{extension structure on} $\mathcal M_{0,3}^{\approx}$
\[
1\longrightarrow\widehat{\mathcal G}_{0,3}^{+}
\longrightarrow\mathcal M_{0,3}^{\approx}
\longrightarrow\mathcal N_{0,3}^{!\approx}\longrightarrow1,
\tag{36}
\]
containing (24) as a subextension of the same quotient by a smaller group.

The subgroup $\mathcal M_{0,3}^{!\approx}(0)$, which is a section of (24), is therefore also a section of (36), hence a canonical isomorphism
\[
\mathcal M_{0,3}^{\approx}
\simeq\underbrace{\unsure{\mathcal N|\mathcal M}_{0,3}^{!\approx}(0)}_{\simeq\mathcal N_{0,3}^{!\approx}}
\cdot\widehat{\mathcal G}_{0,3}^{+}
\tag{37}
\]
($\frac12$-direct product).
\margin{240}

\scanpage{242}

Let us now consider the map
\[
\Phi:\quad
\diagram{\begin{tikzcd}[ampersand replacement=\&,column sep=large]
\mathcal M_{0,3}^{!\approx}\arrow[r]\arrow[d,phantom,sloped,"\simeq"]
\&\mathcal M_{0,3}^{\approx}\arrow[d,phantom,sloped,"\simeq"]\\
\unsure{\mathcal M}_{0,3}^{!\approx}(0)\cdot\widehat\Pi_{0,3}
\&\mathcal M_{0,3}^{!\approx}(0)\cdot\widehat{\mathcal G}_{0,3}^{+}
\end{tikzcd}}
\tag{38}
\]
which is \unsure{given|given} by
\[
\Phi(u\cdot a)=u\widehat\varphi(a)
\tag{39}
\]
($u\in\mathcal M_{0,3}^{!\approx}(0)$, $a\in\widehat\Pi_{0,3}$). It follows from the characterization of $\mathcal M_{0,3}^{!\approx}(0)$ (corollary 3) --- \deleted{\unsure{i.e.}} \inserted{more \uline{particularly}} (from the fact that the elements of $\mathcal M_{0,3}^{!\approx}(0)$ actually commute (not only as outer actions) with $\widehat\varphi$, that $\Phi$ is a group homomorphism. This is of course none other than the homomorphism $\Phi$ considered in \S\,41 (page\,316') \inserted{(up to restriction to the subgroups $\mathcal M_{0,3}^{!}=\unsure{\mathcal E}_{0,3}$ and $\mathcal M_{0,3}$)} --- in loc. cit. this homomorphism could moreover be defined in an effective way --- not only up to $\widehat\Pi_{0,3}$-conjugation or up to $\widehat{\mathcal G}_{0,3}^{+}$-conjugation --- and with a little care we should be able to verify \uline{the equality} of these homomorphisms $\Phi$ from \S\,41 and from here, in an \uline{effective} way.
\margin{241}

\scanpage{243}
\mspage{422}
\uline{Calculation \unsure{with|of}} $x=\ell_\infty$.

I am going to return to the determination of the $u_\xi\in\mathcal M_{0,3}^{!\sim}[0]=\operatorname{Norm}_{\mathcal M_{0,3}^{!\sim}}(L_0)$ which are \uline{normalized} i.e. which commute with $\widehat\varphi$ i.e. satisfy (5) with $\lambda=1$. For $x=\ell_0$ this formula, by virtue of the particular form of $u_\xi$, is tautological
\[
(\widehat\varphi u\ell_0=\widehat\varphi(\ell_0^\mu)=\varepsilon_0^\mu,
\quad u\widehat\varphi(\ell_0)=u(\varepsilon_0)=\varepsilon_0^\mu),
\]
so we \unsure{want|see} that it \unsure{should come|comes} to the same thing to impose it for $x=\ell_1$, or for $x=\ell_\infty$. For $x=\ell_1$, the condition takes the form (10) --- remarkable in that it gives a calculation \unsure{of the|of} ``delicate'' \deleted{the \unsure{variable}}\inserted{\unsure{parameter}} $\beta$ (whose existence \unsure{au priori|a priori}, in terms of $\xi$, rests on foundations which have not yet been established\ldots), in terms of the ``coarse'' parameter $\xi$ --- taking into account the fact that $i$\inserted{$\in\mathbb Z/2\mathbb Z$} itself is determined just as easily in terms of $\xi$.

We are going to give this condition a form rather different in appearance, by writing the commutation relation (5) for $x=\ell_\infty$, instead of $x=\ell_1$. Let us first note that we have the relation
\[
\tau'_\infty(\rho)=\rho^{-1}
\tag{40}
\]
For
\[
\tau'_\infty(\rho)
=\sigma_\infty(\underbrace{\tau_\infty(\rho)}_{\ell_1^{-1}\rho})
=\deleted{\underbrace{\unsure{(\sigma_\infty\ell_1^{-1}\sigma_\infty^{-1})}}_{\unsure{\ell_0^{-1}}}}\inserted{\unsure{\ell_0^{-1}}}
\underbrace{\sigma_\infty(\rho)}_{\underbrace{\sigma_\infty\rho}_{\varepsilon_0}\sigma_\infty^{-1}}
=\varepsilon_0^{-1}\sigma_\infty^{-1}=\rho^{-1},
\]
\margin{242}

\scanpage{244}

it follows that if $u\in\mathcal M_{0,3}^{!\sim}$ satisfies
\[
u(\rho)=\rho^{\unsure{\mu|j}}
\qquad \unsure{\mu|j}\in\mathbb Z/3\mathbb Z\text{ suitable}
\tag{41}
\]
then we necessarily have $j\equiv-1\text{ or }+1\ (3)$ and if \deleted{\unsure{$\alpha$ is} \illegible{1}$(j)=\alpha_3$ is the integer \illegible{1}}
\[
\alpha_3(j)\in\mathbb Z/2\mathbb Z
\quad\text{characterized by}\quad
j\equiv(-1)^{\alpha_3(j)}\ (3)
\tag{42}
\]
we will therefore have $u(\rho)=\rho$ or $u(\rho)=\rho^{-1}$ according as $\alpha_3(j)=0$ or $\alpha_3(j)=1$, and in any case we will have
\[
u(\rho)=\tau'_\infty{}^{\alpha_3}(\rho).
\tag{43}
\]

This being so, let us put $x=\ell_\infty$ in (5):
\[
\widehat\varphi u(\ell_\infty)
=\widehat\varphi(\operatorname{int}(f_\infty^{-1})\ell_\infty^\mu)
=\operatorname{int}(\widehat\varphi(f_\infty^{-1}))\rho^{-\mu}
\]
(since $\widehat\varphi(\ell_\infty)=\rho^{-1}$), and
\[
u(\widehat\varphi(\ell_\infty))
=u(\rho^{-1})=u(\rho)^{-1}
=\deleted{\unsure{(\xi\rho)^{-1}}=\rho^{-1}\xi^{-1}
=\unsure{(\rho^{-1}\xi^{-1}\rho)\rho^{-1}}}
\inserted{\operatorname{int}(\alpha)(\rho^{-1})}
\]
so formula (5) for $\lambda=1$ becomes
\[
\operatorname{int}(\widehat\varphi(f_\infty)\alpha)(\rho^{-1})
=\rho^{-\mu}\qquad\text{i.e.}
\]
\[
\operatorname{int}(\widehat\varphi(f_\infty)\alpha)\rho=\rho^\mu
\]
and by virtue of (43) \deleted{\unsure{this is equivalent} $\unsure{\rho^\mu=\tau'_\infty{}^{\alpha_3}(\rho)}$}\inserted{\unsure{we have}} \deleted{\unsure{that}}
\[
\rho^\mu=\tau'_\infty{}^{\alpha_3=\alpha_3(\mu)}(\rho)
\overset{\unsure{\mathrm{df}|\mathrm{déf}}}{=}
\operatorname{int}(\tau'_\infty{}^{\alpha_3})(\rho)
\]
so the condition becomes
\[
\operatorname{int}(\tau'_\infty{}^{\alpha_3}\widehat\varphi(f_\infty)\alpha)\rho=\rho
\tag{\unsure{44}}
\]
i.e. that the element $\tau'_\infty{}^{\alpha_3}\widehat\varphi(f_\infty)\alpha$ of $\widehat{\mathcal G}_{0,3}=\widehat{\mathcal G}_{0,3}^{+}\cdot\{1,\tau_\infty\}$ commutes with $\rho$.
\margin{243}
